\begin{ouverture}[calcul]{7}{Comparer et additionner des fractions}{Partie 2 · Fractions et nombres décimaux}
{Mêmes parts · Parts compatibles · Sommes}
\objectif{Comparer des fractions de même dénominateur ou de même numérateur.}
\objectif{Rendre deux fractions compatibles quand un dénominateur est multiple de l’autre.}
\objectif{Additionner et soustraire des fractions compatibles.}
\prerequis{Multiples, diviseurs et fractions sur la demi-droite.}
\end{ouverture}
\pagestyle{otMain}\essentiel\fontsize{11}{15}\selectfont
\phantomsection\label{cours-7-1}
\essSection{1}{Comparer des fractions}
\begin{propriete}Même dénominateur : la plus grande a le plus grand numérateur, $\frac79>\frac59$. Même numérateur : la plus grande a le plus petit dénominateur, $\frac35>\frac38$, car les cinquièmes sont plus grands que les huitièmes.\end{propriete}
\begin{savoirfaire}[Changer d’écriture pour comparer]Quand un dénominateur est un multiple de l’autre, convertir : $\frac56$ et $\frac{11}{12}$ se comparent en douzièmes, car $6\times2=12$ : $\frac56=\frac{10}{12}<\frac{11}{12}$.\end{savoirfaire}
\begin{exemple}Comparer à 1 aide souvent : $\frac98>1>\frac78$. Et comparer à $\frac12$ : $\frac38<\frac12$ car $\frac12=\frac48$ ; $\frac7{12}>\frac12$ car $\frac12=\frac6{12}$.\end{exemple}
\begin{attention}$\frac56=\frac{10}{12}$ : multiplier le numérateur et le dénominateur par le même nombre ne change pas le nombre, car les parts deviennent deux fois plus fines mais deux fois plus nombreuses.\end{attention}
\clearpage
\phantomsection\label{cours-7-2}
\essSection{2}{Additionner et soustraire}
\begin{propriete}Avec le même dénominateur, les parts s’ajoutent : $\frac5{12}+\frac4{12}=\frac9{12}$ et $\frac{11}{12}-\frac7{12}=\frac4{12}$.\end{propriete}
\begin{savoirfaire}[Convertir puis réunir]Quand un dénominateur est multiple de l’autre, convertir d’abord : $\frac34+\frac58=\frac68+\frac58=\frac{11}8$, soit $1+\frac38$. Les huitièmes conviennent, car $8$ est un multiple de $4$.\end{savoirfaire}
\begin{exemple}Compléter à l’unité : $\frac7{10}+\frac3{10}=1$. Retirer à une unité : $1-\frac4{15}=\frac{15}{15}-\frac4{15}=\frac{11}{15}$. Retirer à un entier : $2-\frac56=\frac{12}6-\frac56=\frac76$.\end{exemple}
\begin{attention}On n’additionne jamais les dénominateurs : $\frac14+\frac14$ fait $\frac24$, deux quarts, et sûrement pas $\frac28$. Le partage ne change pas quand on réunit des parts.\end{attention}

\clearpage\rubrique{Découverte}
\begin{blocLivre}{ex-7-1}\begin{exercice}[1]{Mêmes parts}
Comparer $\frac79$ et $\frac59$, puis ranger par ordre croissant : $\frac8{15}$ ; $\frac2{15}$ ; $\frac{11}{15}$.
\end{exercice}\end{blocLivre}
\begin{blocLivre}{ex-7-2}\begin{exercice}[1]{Mêmes numérateurs}
Comparer $\frac35$ et $\frac38$, puis $\frac7{10}$ et $\frac7{12}$.
\end{exercice}\end{blocLivre}
\begin{blocLivre}{ex-7-3}\begin{exercice}[1]{Écritures équivalentes}
Compléter : $\frac56=\frac{\square}{12}$ ; $\frac34=\frac{\square}{20}$ ; $\frac{2}{3}=\frac{\square}{60}$.
\end{exercice}\end{blocLivre}
\begin{blocLivre}{ex-7-4}\begin{exercice}[1]{Comparer en convertissant}
Comparer $\frac56$ et $\frac{11}{12}$, puis $\frac23$ et $\frac7{12}$.
\end{exercice}\end{blocLivre}
\begin{blocLivre}{ex-7-5}\begin{exercice}[1]{Autour de un demi}
Comparer chaque fraction à $\frac12$ : $\frac38$ ; $\frac7{12}$ ; $\frac5{10}$.
\end{exercice}\end{blocLivre}

\rubrique{Entraînement}
\begin{blocLivre}{ex-7-6}\begin{exercice}[2]{Sommes directes}
Calculer $\frac5{12}+\frac4{12}$, puis $\frac{11}{15}-\frac4{15}$.
\end{exercice}\end{blocLivre}
\begin{blocLivre}{ex-7-7}\begin{exercice}[2]{Convertir puis additionner}
Calculer $\frac34+\frac58$, puis $\frac23+\frac16$.
\end{exercice}\end{blocLivre}
\begin{blocLivre}{ex-7-8}\begin{exercice}[2]{Retirer à un entier}
Calculer $1-\frac4{15}$, puis $2-\frac56$.
\end{exercice}\end{blocLivre}
\begin{blocLivre}{ex-7-9}\begin{exercice}[2]{Compléter}
Compléter : $\frac7{10}+\square=1$ ; $\frac38+\square=\frac78$ ; $\square-\frac5{12}=\frac1{12}$.
\end{exercice}\end{blocLivre}
\begin{blocLivre}{ex-7-10}\begin{exercice}[2]{L’erreur classique}
Mehdi écrit : $\frac14+\frac14=\frac28$. Expliquer son erreur avec une image de parts, puis corriger.
\end{exercice}\end{blocLivre}

\rubrique{Problèmes}
\begin{blocLivre}{ex-7-11}\begin{exercice}[3]{La tarte partagée}
Léa mange $\frac38$ d’une tarte, Max en mange $\frac14$.\begin{questions}
\item Quelle fraction de la tarte ont-ils mangée ensemble ?
\item Quelle fraction reste-t-il ?
\end{questions}
\end{exercice}\end{blocLivre}
\begin{blocLivre}{ex-7-12}\begin{exercice}[3]{Le trajet en trois temps}
Un trajet est fait aux $\frac7{12}$ en bus, à $\frac14$ en tram, et le reste à pied. Quelle fraction du trajet est faite à pied ?
\end{exercice}\end{blocLivre}
\begin{blocLivre}{ex-7-13}\begin{exercice}[3]{Les bouteilles d’eau}
Un sportif boit $\frac34$ de litre le matin et $\frac58$ de litre l’après-midi. A-t-il bu plus ou moins d’un litre ? De combien ?
\end{exercice}\end{blocLivre}
\begin{blocLivre}{ex-7-14}\begin{exercice}[3]{Le classement des lectures}
Trois élèves ont lu le même livre. Ana en est aux $\frac7{10}$, Bilal aux $\frac35$, Chloé aux $\frac34$. Qui est le plus avancé ? Ranger les trois avancées.
\end{exercice}\end{blocLivre}
\begin{blocLivre}{ex-7-15}\begin{exercice}[3]{Un pas vers l’abstrait}
Sans convertir, expliquer pourquoi $\frac5{13}+\frac8{13}=1$, puis calculer $\frac{37}{60}+\frac{23}{60}$.
\end{exercice}\end{blocLivre}

\clearpage
